\section{Architecture générale}
\label{sec:archi}

Le firmware est organisé en couches à responsabilité unique, dont la règle de dépendance est stricte : une couche ne peut dépendre que de couches situées en dessous d'elle, jamais l'inverse~\cite{martin_clean_architecture}. Cette règle est matérialisée par le découpage en composants ESP-IDF, chaque composant déclarant explicitement ses dépendances dans son \texttt{CMakeLists.txt}~\cite{ESP32_IDF_guide}. La figure~\ref{Architecture_logicielle.drawio} présente cette organisation.

\fig[H, width=0.9\textwidth]{Architecture logicielle en couches. Les traits pleins représentent les dépendances de compilation. Diagramme créé avec l'aide de Claude Opus 4.8.}{Architecture_logicielle.drawio}

La couche \texttt{bsp} est la seule à connaître le brochage de la carte, regroupé dans \texttt{board\_pins.h}. Elle initialise les bus partagés (ADC, I\textsuperscript{2}C, I\textsuperscript{2}S, SPI) et en transmet les handles aux couches supérieures, si bien qu'aucun pilote ne configure lui-même son contrôleur. La couche \texttt{drivers} contient un pilote fin par circuit. La couche \texttt{services} expose des capacités indépendantes du matériel : stockage, pipeline audio, Bluetooth, entrées utilisateur, énergie et préférences persistantes. La couche \texttt{ui} fournit le cadre graphique générique (framebuffer, navigation, menus) et la couche \texttt{screens} les écrans concrets. Le module \texttt{main} assemble le tout, il initialise chaque couche dans l'ordre de ses dépendances, crée les tâches, puis se retire du fonctionnement courant.

La règle de dépendance porte sur le sens et non sur l'adjacence. Une couche peut court-circuiter celle qui la précède immédiatement lorsqu'elle n'en a pas l'usage. La sortie audio filaire appartient ainsi aux \texttt{services} mais pilote directement le bus I\textsuperscript{2}S du \texttt{bsp}, aucun pilote n'ayant à s'intercaler. L'inverse reste interdit, aucun module de bas niveau ne connaît les couches situées au-dessus de lui.

\subsection{Modèle d'exécution}

Le firmware repose sur quatre tâches FreeRTOS permanentes et une conception entièrement événementielle, résumées au tableau~\ref{tab:taches}. La tâche audio, prioritaire et épinglée sur le cœur~1, exécute la boucle de décodage et d'écriture vers la sortie active. La tâche d'interface, également sur le cœur~1, réalise le rendu et la navigation. Ce placement les isole du cœur~0, que la pile Bluetooth sature pendant le streaming~\cite{ESP32_IDF_guide}. La tâche d'entrée traduit les interruptions des boutons en événements. La tâche de maintenance regroupe enfin les tâches périodiques sur une base de 100~ms : lecture du potentiomètre de volume, détection des insertions et des retraits de la carte SD, détection de la perte du lien Bluetooth, télémétrie batterie, échantillonnage d'une éventuelle campagne d'autonomie et conditions d'arrêt automatique.

\begin{table}[H]
    \centering
    \caption{Tâches FreeRTOS permanentes du firmware. Les piles indiquées sont les tailles allouées à la création.}
    \label{tab:taches}
    \footnotesize
    \begin{tabular}{lcccl}
        \toprule
        Tâche       & Cœur  & Priorité & Pile   & Rôle                                  \\
        \midrule
        audio       & 1     & 22       & 4\,kio & décodage et écriture vers la sortie   \\
        ui          & 1     & 5        & 8\,kio & rendu, navigation, verrouillage écran \\
        input       & libre & 5        & 3\,kio & lecture des boutons et mise en file   \\
        maintenance & libre & 3        & 3\,kio & tâches périodiques (100~ms)           \\
        \bottomrule
    \end{tabular}
\end{table}

Aucune tâche n'utilise d'attente active. Chaque tâche bloque sur sa file d'événements au repos, libérant ainsi le processeur sans consommer de cycles~\cite{barry_mastering_freertos}. Un appui bouton génère une interruption sur l'expanseur qui débloque simplement la tâche d'entrée, le contexte d'interruption interdisant les transactions I\textsuperscript{2}C nécessaires à l'identification~\cite{diodes_inc_PI4IOE5V9554}. Cette tâche lit l'état des boutons, génère les événements correspondants et les transmet à la tâche d'interface. De même, la tâche audio reste bloquée en l'absence de commande de lecture.

Seule la tâche de maintenance utilise la scrutation pour surveiller les grandeurs sans interruption comme le potentiomètre ou la batterie. Sa période de 100~ms résulte d'un compromis entre réactivité du réglage du volume et consommation énergétique. Ce modèle événementiel préserve l'autonomie car les tâches bloquées ne consomment aucun cycle, permettant au processeur d'entrer en veille profonde entre deux réveils~\cite{ESP32_IDF_guide}.

À ces quatre tâches permanentes s'ajoutent des tâches éphémères, créées à la demande pour les opérations longues et bloquantes puis détruites à leur terme : indexation de la carte SD, réinitialisation ou extinction de la pile Bluetooth, commutation du multiplexeur USB. Ce procédé évite qu'une opération de plusieurs centaines de millisecondes ne bloque la tâche d'interface ou la tâche de maintenance, sans immobiliser en permanence la pile que cette opération réclame.

Le partage des ressources obéit lui aussi à une propriété unique. Chaque périphérique SPI dispose de son propre hôte matériel, l'écran sur SPI2 et la carte SD sur SPI3, ce qui permet à la tâche audio de lire la carte pendant que la tâche d'interface rafraîchit l'écran sans aucune exclusion mutuelle~\cite{ESP32_Technical_Reference_Manual}. Le bus I\textsuperscript{2}C, partagé entre le DAC, la jauge et l'expanseur, est possédé par le \texttt{bsp} qui en distribue les handles.

Le placement mémoire répond à une contrainte matérielle : le contrôleur DMA de l'ESP32 ne peut adresser la PSRAM~\cite{ESP32_WROVER-E_Datasheet}. Les tampons utilisés par les périphériques doivent donc résider en RAM interne, une ressource limitée fortement sollicitée par la pile Bluetooth~\cite{ESP32_IDF_guide}. À l'inverse, les données accessibles uniquement par le processeur, comme le framebuffer ou l'index des pistes, sont placées en PSRAM pour préserver la mémoire interne.

Ce modèle d'exécution a été validé par un test d'endurance de 3~h~34 dont les résultats figurent à la section~\ref{sec:res_execution}.
